如果你决定读这本书，你或许已经对 Rust 如何在职业或个人项目中为你带来帮助有所了解。也许你用过用 Rust 编写的某些软件，想更深入地了解这门语言。又或者你只是出于挑战自我，想学一门新的编程语言。甚至可能是某位好心的同事或朋友告诉你，Rust 是“C++ 杀手”，或者学会它能帮你写出更好的 C++。
无论出于什么原因，我们都不妨在开篇花点时间回答一个问题：程序员最初究竟为何应该关注 Rust。我们将考察 Rust 从何而来，看一些使用它的软件实例，了解它在不断扩张的编程语言领域中处于什么位置，并探索如何借助 Rust 语言、工具链和社区来帮助我们打造出色的软件。
本章我们将涵盖以下主要主题：
\begin{itemize}
  \item Rust 简史
  \item Rust 如何改进软件
  \item crate 与 Rust 生态
\begin{myNotic}{注意}
\textbf{购买本书即享免费 PDF 副本 + 专属额外内容}

购买本书即可获得本书的无 DRM PDF 副本、Packt+ 书库 7 天试用（无需信用卡），以及额外的专属内容。请参阅\emph{前言}中的\emph{随书附赠的免费权益}一节，立即解锁它们，让学习效果最大化。
\end{myNotic}
\end{itemize}
\section{Rust 简史}
\label{idx:idx_cea019cd}
考察一门编程语言时，了解它被创造时所处的历史背景会很有帮助。它当初是为了解决什么问题而诞生？它又选择了怎样的取舍？它最吸引哪一类人？为了回答这些问题，我们将跟随 Rust 的历程，看它如何从一个人的个人项目，发展成如今驱动着数百万设备上的 Web 浏览器和操作系统的语言。
\subsection{Mozilla 的早期起步}
\label{idx:idx_b6a784b3}
Rust 于 2006 年由 Graydon Hoare 作为个人项目启动，他是就职于 Mozilla Research 的工程师。这门新语言从一开始就被设计成实用之作，而不是一个永远默默无闻的研究项目。事实上，他明确瞄准的是“受挫的 C++ 开发者”，切入的是编程语言版图中传统上由 C 和 C++ 占据的位置。
C 和 C++ 与其他编程语言的区别主要体现在两个领域：它们赋予开发者对内存的控制程度，以及它们将语言强加的运行时开销降到最低的做法。许多语言至少会在这两个领域之一做出妥协。这并不一定使它们成为糟糕的语言，只是让它们各自适用于不同的用途。像 C 和 C++ 这类语言所选择的取舍，对许多软件项目而言依然合理，尤其是那些无法容忍任何性能或内存开销的项目。
正如我们对任何瞄准受挫 C++ 开发者的语言所期待的那样，Rust 的设计做出了与 C 和 C++ 类似的取舍：在最小化运行时开销的同时最大化控制力。最小化运行时开销是一个容易的选择，但最大化对内存的控制力却有一些棘手之处，Rust 的设计需要重新审视。即便是最优秀的程序员也会犯错。而当程序员在内存上犯错时，后果可能从时隐时现的怪异行为，一直到可被利用的安全漏洞——后者让黑客牟利，却牺牲了我们的用户。
现在我们就能开始想象，为什么 Graydon Hoare 和他身边的 C++ 程序员会感到受挫了。他们想少犯错，但手头可用的工具并没有真正提供他们所需的帮助。Rust 是他们赠予我们的礼物——赠予我们这些不断犯下同样错误、需要一点帮助的 C++ 程序员。Rust 的思路，是让我们把代码划分为“安全”和“不安全”两类，从而给予我们这份帮助。我们的 Rust 代码大部分都会是“安全”的，在一套编译期限制下运行，这些限制协同作用，把我们从许多具体的内存处理缺陷中解放出来。
对于需要额外控制的那一小部分代码，我们可以使用“不安全”的 Rust 代码，并依靠测试和其他验证方法来证明它工作正确。如果我们选择得当，这会大幅缩小我们必须煞费苦心地验证的代码范围。对大型软件项目而言，曾经艰巨无比的任务突然变得可以完成。
\label{idx:idx_750b6958}
\label{idx:idx_aed0c55d}
当 Mozilla 看到 Hoare 一直在开发的 Rust 原型后，他们很快意识到它所要解决的问题，并决定让这门新语言在他们最熟悉的软件——Web 浏览器——上发挥作用。Mozilla 启动了一个名为 \textbf{Servo} 的项目，其主要目的是从零开始、用 Rust 打造一个全新且安全可靠的浏览器引擎。Servo 项目为 Rust 的早期开发提供了大量资金，并把它从一个人的个人项目，一举推成拥挤的编程语言舞台上一位强有力的竞争者。但当时大多数科技公司还没有准备好把它们的软件项目迁往 Rust。
\subsection{让语言趋于稳定}
\label{idx:idx_a0d14c34}
在语言发展早期编写的代码，在今天的 Rust 程序员看来可能已经面目全非。由于开发者要理清他们想要创造的语法和标准库，语言的某些部分变化得很快。这样的变化速度对于语言演变成今天我们所熟知的样子是必要的，但这也意味着，当时任何想使用 Rust 的人都必须随着每个新版本的发布更新、甚至重写自己的代码。这样的局面，无法激起更广泛的软件行业对这门语言的认真、长期的兴趣。
2015 年 5 月，这种局面得到解决，Rust 达到了稳定的 1.\allowbreak{}0 版本。经历了如此多的变动之后，开发者们终于确信，这门语言的功能集不再需要随着每次发布而大幅变动。所谓“稳定”并不意味着这门语言离完工有多近，但 1.\allowbreak{}0 版本附带了一项稳定性保证。开发者承诺，从此以后编写的任何 Rust 程序，在使用未来任何版本的 Rust 时都无需修改即可继续运行。但是，一门带有这种稳定性保证的语言，岂不是会无限膨胀下去？尤其是在那些已经存在很久的语言里，无法废弃某些功能，对用户和语言设计者来说都是一根扎在肉里的刺。C++ 应对这一问题的办法，是选择性地在一些被认为影响较小的位置打破兼容性，但即便如此，语言的演进仍受到必须支持长达数十年遗留代码这一要求的拖累。
\label{idx:idx_7d1e474a}
\label{idx:idx_5fdaed9b}
看到 C++ 和其他长寿命语言所遭遇的问题后，Rust 的设计者从一开始就知道要为这种情况做好规划。他们设计出一个名为 \textbf{版本（edition）} 的概念，用来对语言做出改动。新的 edition 可以对语言做出重大改动，包括新增保留关键字、废弃并移除功能，以及更改语言语义。Rust 大约每隔三年发布一个新的 edition。2018 edition 和 2021 edition 各自都为这门语言带来了大量有趣的增删。虽然大多数 Rust 程序员会选择把代码迁移到最新的 edition，以享受新功能带来的好处，但编译器会无限期地支持较旧的 edition，从而为任何认为把项目推进到新 edition 风险太大或太耗时的人提供兼容性保证。
稳定的 1.\allowbreak{}0 版本已经发布，语言未来如何演进也有了一套可靠的方案，Rust 终于为 Mozilla 之外的公司在其自有项目中试用这门语言做好了准备。果不其然，在稳定之后的数年里，这门语言在更广泛的科技行业中开始了缓慢而稳步的流行度增长。
\subsection{拓展新领域}
\label{idx:idx_9fc025cf}
Rust 最早的几场重大胜利之一，就来自 Mozilla 的 Servo 项目。在这个新浏览器项目开发数年后，Mozilla 决定开始用为 Servo 开发的 Rust 组件替换 Firefox 浏览器中的部分 C++ 代码。从 2017 年底开始，Firefox 向最终用户发布的版本开始启用 Rust 代码。尽管安全是 Servo 项目的主要关注点之一，但运行速度的提升也随之而来。这个被命名为“Quantum”的 Firefox 版本，凭借在 Servo 中用 Rust 实现的 CPU 并行化改进，在渲染速度上实现了显著提升。
早在 2018 年，Amazon 就开始使用 Rust 来改进其 AWS 云服务。他们对 Rust 的公开支持，始于 Amazon 宣布 Firecracker VM，它为 AWS 的部分无服务器产品提供支持。在 2020 年的一篇博客文章中，Amazon 称赞 Rust 对性能与可靠性的关注，同时透露若干其他 AWS 服务，包括其云存储和计算产品，如今也在以某种形式使用 Rust。对 Rust 的这种依赖，促使 Amazon 聘请了若干关键的 Rust 贡献者，以确保这门语言能够持续满足他们的需求，也能满足更广泛社区的需求。
同样，Rust 早在 2019 年就引起了微软的注意。负责处理微软产品安全问题管理的微软安全响应中心表示，他们处理的最严重安全问题中有 70\% 是由内存安全问题引起的……而这正是 Rust 志在防范的东西。微软发布了一些有助于使用 Windows API 的开源 Rust 代码，但此时还不清楚他们究竟打算如何在自家产品中使用 Rust。
与此同时，Google 正悄然在其 Android 操作系统和 Chromium 浏览器中拥抱 Rust。Google 的投入之大令人惊讶；仅 2022 年一年，他们就统计出内部有超过 1,\allowbreak{}000 名个人贡献者参与了其 Rust 项目。Google 还对其 Rust 程序员应用了内部的分析流程，发现其中三分之二的人认为，自己能在两个月或更短时间内具备为 Rust 代码库做贡献的信心。大多数面向 Rust 初学者的资料都强调这是一门学习曲线陡峭的语言，但 Google 的经验表明，对有经验的工程师来说，它或许并非一道很高的门槛。
Rust 最令人意外的采用者之一，竟是历史上用 C 和汇编语言编写的 Linux 操作系统。2007 年，当贡献者建议允许把 C++ 引入内核时，首席维护者 Linus Torvald 对 C++ 的激烈抨击一度广为人知。因此，当他对把 Rust 作为内核开发语言的努力表现出谨慎的支持时，许多人都大为震惊。经过一小群维护者数月的细致工作，Rust 于 2022 年底被合并进内核，完成了 C++ 未能做到的事。
\label{idx:idx_5d8e82a0}
2023 年 3 月，微软又带来一个意外消息，终于让人看清了他们一直如何在幕后使用 Rust。微软从 2020 年起就在试验用 Rust 重写 Windows 操作系统的部分内容，如今已准备开始向终端用户交付。2023 年 5 月，Windows 11 的首批 insider 构建版本向用户交付了 Rust 代码，这让 Windows 成为继 Android 和 Linux 之后第三个随用 Rust 编写的内核代码一同发布的重大操作系统。与 Linux 一样，Windows 内核中的 Rust 代码量仍然很小，但微软工程师已表示，他们打算在初步试验之外进一步扩大使用范围。
如今，按用户数量排名前五的操作系统中有三个、三大 Web 浏览器后端中的两个，以及云厂商的若干产品，都被认为用到了 Rust。Rust 吸引了这些领域以及其他更多领域中感到不满的开发者，而这些领域历来多由 C 和 C++ 之类的语言主导。尽管大多数系统级代码仍用 C 或 C++ 编写，但 Rust 在科技行业已拥有清晰且快速增长的立足点，这会让它在未来许多年里始终保持重要地位。
\section{Rust 如何改进软件}
\label{idx:idx_a2e0f75c}
要让我们选择 Rust，它必须提供一组极具吸引力的特性，能相对现状带来实质性的改进。本节中，我们将考察 Rust 的一些方面和特性，正是它们让 Rust 成为构建可靠、高性能软件的绝佳选择。
\subsection{高层与底层}
\label{idx:idx_e8949b99}
我想先分享一个我最喜欢的、适合 Rust 初学者的例子：读取一个文件并生成哈希值。在 C++ 中解决这个问题，大概要从大量搜索开始，去找一个许可合适、既实现了所需算法、又恰好与我们所用构建系统兼容的库。C++ 中有若干工具能帮助解决这个问题，但它们没有一个是每套 C++ 工具链都会默认安装的。
我之所以喜欢用这个问题，只是为了说明 Rust 开箱即用、易于上手。获取依赖轻而易举，而且 Rust 里的简单程序看起来确实很简单。不要被陌生的语法或程序的具体细节吓到——这些我们在阅读本书的过程中都会学到。相反，我希望你关注整个过程有多简洁。要新建一个 Rust 项目，我可以在终端或 shell 应用中创建一个骨架项目，再敲几秒钟键盘，加一个合适的库来替我完成哈希计算。本例使用的是 Rust 1.\allowbreak{}82.\allowbreak{}0，但任何更新的 Rust 版本行为都完全相同：
\begin{codeboxdark}
$ cargo init hash
$ cd hash
$ cargo add sha256
\end{codeboxdark}
程序本身在 Cargo 创建的 \texttt{src/main.rs} 文件中只有寥寥几行：
\begin{codebox}[rust]
use sha256::digest;
fn main() -> Result<(), std::io::Error>    {
    for argument in std::env::args().skip(1) {
        let file = std::fs::read(&argument)?;
        let hash = digest(file);
        println!("{}: {}", argument, hash);
    }
    Ok(())
}
\end{codebox}
之后，编译并运行这个程序可以一步完成：
\begin{codeboxdark}
$ cargo run --- src/main.rs
\end{codeboxdark}
这里把包含我们代码的 \texttt{src/main.rs} 文件作为参数传给程序，程序则会打印出该文件的哈希值：
\begin{codeboxdark}
src/main.rs: 327268ae21b5c28adea2e290af7f35bca3baa1089fa7bd4a22193743196f7eb8
\end{codeboxdark}
\label{idx:idx_18628be7}
在审视这样的程序时，我喜欢引用 Alan Perlis 关于高级与低级编程语言的观点，它出自他那本经久不衰的\emph{Epigrams on Programming}：
“\emph{当程序需要关注无关紧要的东西时，这门编程语言就是低级的}。”
虽然前面那个 Rust 程序可能并不完美，但按 Perlis 的标准，它显然不算低级——原因很简单：程序员能够轻松、安全地忽略掉大量内容。内存分配和错误处理被隐藏在直白的抽象之后，甚至连变量的类型都被省略了。由于我完全用安全 Rust 编写，还能确信程序不会崩溃，也不会触发任何意外行为。
这套工具链用起来同样毫不费力。创建新项目或添加依赖，我几乎不用多想。我确实得去查一下要添加的依赖的名称和文档，但正如接下来会看到的，这通常也毫无痛苦。由于这门语言具有高级语言的特质，并自带静态安全保证，工具箱里有 Rust 的程序员能快速写出扎实、性能良好的程序。
\label{idx:idx_a04aa939}
然而，Rust 也允许我们深入底层编程，深度足以解决问题即可。如果内存或文件 I/\allowbreak{}O 的安全抽象给不了我们足够的控制力，我们可以自行创建内存分配或读取缓冲策略。如果 Rust 的安全检查妨碍了我们进行某项优化，只要愿意承担风险，我们就可以把它们关掉。
但默认情况下，对于简单的任务，Rust 往往只暴露简单的 API。事实上，总体而言，Rust 的一大优势就是默认做出出色的选择。
\subsection{默认即安全}
\label{idx:idx_8161d7e1}
微软、Google、Mozilla 以及许多其他公司都在投资 Rust，因为他们都从自己的软件产品中注意到了同一个问题。这些公司内部的不同团队研究了统计数据，发现其产品中最严重的 bug 和漏洞约有 70\% 源于内存安全问题，比如缓冲区溢出，或在内存被释放后继续使用它。
Rust 默认的“安全”配置可以防止由内存误用引发的最常见 bug。虽然已有若干内存安全的语言，但其中许多都做了妥协，以垃圾回收的形式引入运行时开销。Rust 的做法，正如我们将在\emph{第 3 章}中看到的，是通过静态追踪内存在整个程序中的所有权与生命周期流转，从而不产生显著的运行时开销。
Rust 还能帮助我们编写健壮的多线程代码。Rust 的设计者认识到，那些保护程序员免于内存 bug 的编译期规则，同样可以扩展用来保护用户免遭一种常见的多线程问题——\textbf{数据竞争}。正如 Mozilla 的 Servo 团队所发现的，这能带来显著的速度提升：在其他语言中被认为太难或太棘手、难以做对的代码里，它创造出了并行的机会。归根结底，Rust 默认的安全配置保护程序员和用户免于许多代价高昂的错误，而 C++ 开发者则直接暴露在这些错误之下。诚然，C++ 社区围绕标准化和工具链做了巨大努力来消灭这些 bug，但捕获它们的过程依然繁重而昂贵。
\label{idx:idx_968db591}
要在 C++ 项目中达到 Rust 那样的保证，团队必须使用语言特性的某个特定子集、某份特定的风格指南、确保一切合规的静态检查器，以及特殊的运行时配置，例如 LLVM 的未定义行为检测器（undefined behavior sanitizer）、地址检测器（address sanitizer）或线程检测器（thread sanitizer）。而在 Rust 里，所需的一切只是编译并运行程序。
\subsection{性能与内存}
\label{idx:idx_a1c5f2d1}
能够编写开销极小的程序，这一点 C++ 做对了，而任何缺少这种能力的语言都不可能成为 C++ 的强劲对手。令人高兴的是，在性能和内存使用方面，Rust 提供了与 C++ 相似的特性组合，这意味着它开箱即可在这两方面都表现不俗。
与 C++ 的 Clang 编译器一样，Rust 构建在 LLVM 项目的编译器架构之上。这为 Rust 提供了世界级的优化编译器后端，开箱即支持多种不同的 CPU 架构。许多 C++ 程序员熟悉的若干其他语言级优化，比如移动语义和\textbf{单指令多数据}（\textbf{SIMD}）内建函数，也都是 Rust 的特色，我们将在后面的章节中看到。
总体而言，Rust 的内存效率也非常高。Rust 默认没有垃圾回收开销，而是用作用域来管理分配和释放。不过，Rust 在内存分配策略上可以相当灵活。如果另一种分配器能为特定用例带来更好的性能，默认分配器可以被全局替换掉。此外，面向微控制器的嵌入式 Rust 程序通常运行在一种移除堆分配的模式下。这时可以用不使用堆的同类容器的固定大小版本，或其他任何符合程序需求的方式来替代。社区创建的开源库集合中体现了许多不同的分配方法。
\subsection{现代的语言与工具}
\label{idx:idx_ba50838d}
我们已经看到，使用 Rust 的工具链来创建项目、添加依赖有多容易。这不过是 Rust 仅凭“更新”这一点就做对的众多事情之一。这种规律在软件界反复出现。新项目往往能比竞争对手跑得更快，因为它可以吸取老项目遭遇的错误和挑战的教训，并从零开始进行设计以避开它们。
\label{idx:idx_010183fe}
\label{idx:idx_480d1e46}
Rust 有机会向 C++ 所推崇的结构化与面向对象风格之外的语言学习。其语言设计中有一条虽小却清晰可辨的函数式编程脉络。最初的 Rust 编译器是用一门优美的函数式语言\textbf{OCaml}写成的，所以这种影响对我们来说应该并不意外。Rust 中来自函数式编程的最明显特性称作\textbf{和类型}（sum type），在某些圈子里也叫\emph{代数数据类型}（algebraic data types），但对大多数 C++ 程序员来说，最好把它理解为\emph{带标签联合}（tagged unions）。正如我们将在\emph{第 2 章}中看到的，带标签联合能实现方便而强大的抽象，并且是 Rust 错误处理机制的核心。
\label{idx:idx_020dc836}
\label{idx:idx_b2e77528}
Rust 还直接内建了其他一些有趣的类型。\textbf{范围}（range）表示一串连续的整数，可以是有限的，也可以是无限的。\textbf{切片}表示对一串连续值的引用，由指针和计数一起存储，就像 C++20 的
\texttt{std::span}。这使得用切片进行裁剪或调整大小的操作只需 \texttt{O(1)} 的代价。迭代器在 Rust 中也
是内建的，不过它们与 C++ 的迭代器关系不大。相反，它们为高效地操作对象列表提供了函数式接口，这可以成为表达许多循环结构的强大方式。
自 C++ 最初被设计出来后的这些年里，软件测试的重要性显著提升。Rust 的设计者认识到了这一点，把测试运行器加入了语言和构建工具链中。同样，代码格式化和文档生成功能也是任何一次全新安装 Rust 时默认附带的。对于任何一门严肃的新编程语言来说，这些特性几乎已是默认预期，但遗憾的是，许多最流行的 C++ 工具链默认并不提供这些功能。
我们已经简要见识过 Cargo，也就是 Rust 的构建与包管理方案。接下来，我们将探讨这个工具以及更广阔的 Rust 生态如何为我们的程序带来显著提升。
\section{crate 与 Rust 生态}
\label{idx:idx_a44cb325}
通常，较新的语言能否生存，不仅取决于其语言设计的实力，还取决于通过其标准库和社区所提供的库的实力。如果每个上手这门语言的人都发现自己得亲手实现一个 JSON 解析器或 socket 库，那么哪怕这世界上设计得最有趣的语言也无法获得任何关注。由于围绕 Rust 的普遍热情，以及它最初为创建 Servo 浏览器而产生的开发需求，Rust 对最常见的库和技术都有很好的库覆盖。对我们大多数人来说，这意味着可以少花时间到处寻找合适的依赖，多花时间写自己的程序。
\subsection{crate 的解剖}
\label{idx:idx_aeacb72d}
\label{idx:idx_a3f4f844}
\label{idx:idx_6732412a}
Rust 把程序和库依赖都称作\textbf{crate}。这一术语似乎没有什么确凿的技术理由，但把它们理解为“一箱箱代码”（crates of code）或许说得通。毕竟在 Rust 中，crate 定义了编译单元。C++ 把每个 \texttt{.cpp} 源文件定义为编译单元（通常称为\textbf{翻译单元}），而在 Rust 中，编译单元包含构成单个 crate 的一个或多个源文件。
\label{idx:idx_5f04d853}
crate 既可以是库，也可以是单独的程序，后者称为\textbf{二进制 crate}。每个 crate 的根目录下都有一个 \texttt{Cargo.toml} 文件，列出它所依赖的依赖项。库 crate 的顶层文件是 \texttt{src/lib.rs}，而二进制 crate 则使用 \texttt{src/main.rs}。库 crate 还可以随之附带一个或多个二进制程序，它们可以作为单独的文件放在 crate 下的 \texttt{src/bin/} 文件夹中，每个二进制程序的名称就是对应源文件的名称。crate 的 \texttt{src/} 文件夹下的每个额外源文件都代表一个\textbf{模块}，可供 crate 中的其他源文件导入。函数和数据类型既可以是某个模块私有的，也可以暴露给 crate 的其他部分。
\label{idx:idx_b0294c9f}
在源代码布局方面，这些文件和路径体现了 Rust 的“约定优于配置”理念。因为编译单元是 crate，而且 crate 内的源文件布局非常严格，所以任何 crate 都能用一套简单的规则来编译。C++ 有多个构建系统，但每个都需要处理预处理器头文件和宏定义的复杂性，同时还要应付那些在源文件布局或包含方式上没有统一约定的库。几乎每个 C++ 项目都要用不同的方式来处理。Rust 坚持不懈地避免这些问题，这正是它的构建与依赖管理工具 Cargo 如此受推崇的主要原因之一。
\subsection{添加对社区 crate 的依赖}
\label{idx:idx_ea139784}
库 crate 和二进制 crate 都可以上传到 Rust 社区代码的主仓库 \href{https://crates.io/}{\texttt{https:/\allowbreak{}/\allowbreak{}crates.\allowbreak{}io/\allowbreak{}}}。这个网站提供了一份可搜索的全部社区 crate 列表，它们都能轻松导入到 Rust 程序中。正如前面哈希示例中所见，添加依赖可以简单到在 crate 的根目录下使用 \texttt{cargo add} 命令。在构建时，Cargo 会对照所需的 crate 版本检查本地缓存，并下载每个需要编译的 crate 的源代码。
\label{idx:idx_2daa04fa}
\begin{myNotic}{注意}
\textbf{关于代码许可}

始终要检查你导入项目的代码的软件许可证，确保它与你的软件项目兼容。Rust 社区中的许多人都选择以 Apache 和 MIT 许可证对其代码进行双许可，这样代码便可用于闭源项目。但情况并非一概如此。在 \href{https://crates.io}{\texttt{https:/\allowbreak{}/\allowbreak{}crates.\allowbreak{}io}} 上有几个不错的工具可以检查程序依赖项的许可证。
\end{myNotic}
\label{idx:idx_62b9b13f}
这种便利程度可能是 C++ 程序员不熟悉、甚至会心存疑虑的。我们为什么要在构建时下载代码呢？这难道不会把我们绑死在可能随时宕机的外部基础设施上，导致我们无法构建项目吗？简而言之，是的，这确实是一种风险。但在 C++ 生态之外，这是一种常见的风险——JavaScript 和 Python 等语言都有类似的包仓库。如果这是你的组织所担忧的问题，或者只是你作为程序员不愿承担的风险，Rust 有办法绕开它，我们将在\emph{第 4 章}中探讨。
\section{小结}
我们了解到了 Rust 最初是如何为吸引感到沮丧的 C++ 程序员而开发的，也注意到若干重要软件背后的团队正逐步将工作迁移到 Rust 上。我们探讨了这门语言的一些安全保证和语言特性，以及它出色的工具链和由社区打造的强大的开源代码库。
综合来看，这些理由足以让人信服地选择 Rust，但投身这门语言最好的理由远胜于此——用起来很愉快。得益于其强大的安全保证，我们无需再担心在其他语言中编码时必须操心的许多事情。编写或审查 Rust 代码时，会有一种奇妙的自由感。你只需专注于程序本身的真正逻辑，而不用时刻保持警惕，去避免崩溃和内存安全问题这类严重失误。
在\emph{第 2 章}中，我们将更深入地考察 Rust 的语法，并学习如何更有效地使用它的主要语言特性。
\needspace{26\baselineskip}
\begin{samepage}
\subsection{获取本书的 PDF 版本及更多内容}
扫描二维码（或访问 \texttt{packtpub.\allowbreak{}com/\allowbreak{}unlock}）。按书名搜索本书，确认版本，然后按照页面上的步骤操作。
\myGraphic{0.32}{images/qr_unlock.png}{}
\myGraphic{0.20}{images/unlock_btn.png}{}
\emph{注意：请把发票放在手边。直接从 Packt 购买则无需发票。}
\myGraphic{0.32}{images/qr_unlock.png}{}
\myGraphic{0.20}{images/unlock_btn.png}{}
\end{samepage}
